\pagestyle{plain}
\chapter*{Resumen}
\addcontentsline{toc}{chapter}{Resumen}

Este Trabajo de Fin de Grado desarrolla la primera fase de un laboratorio
virtual para estudiar paradigmas de toma de decisiones humanas mediante agentes
inteligentes. La contribución principal es un pipeline que transforma una
descripción en lenguaje natural en modelos de decisión ejecutables en Python.
Para ello se divide el problema en etapas especializadas: investigación
científica, formalización matemática, adaptación al entorno de simulación y
generación de código validado mediante pruebas.

El sistema se apoya en un enrutador determinista con revisión humana entre
etapas, almacenamiento de artefactos en MinIO, metadatos y ciclo de vida en
PostgreSQL, memoria semántica en Qdrant y un grafo de conocimiento en Neo4j. La
arquitectura de memoria almacena solo artefactos aceptados, combina recuperación
densa, BM25 y recorrido de grafo, y aplica mecanismos de corrección para evitar
que conocimiento irrelevante o desactualizado condicione nuevas ejecuciones.

El resultado de la fase es un conjunto de modelos compatibles con un contrato
simple de simulación: \texttt{decide}, \texttt{update} y \texttt{get\_state}.
Esta separación permite que la segunda fase del proyecto cargue dinámicamente
los modelos generados sin depender de una jerarquía rígida de clases.

La validación combina pruebas unitarias, integración continua, evals internas y
evals con corpus cerrado de papers. Estas pruebas muestran que el pipeline
puede completar ejecuciones largas y producir evidencia auditable. Aun así, la
validación técnica no sustituye la revisión científica por expertos del dominio.

\vspace{0.5cm}
\noindent{\bf Palabras clave:} agentes inteligentes, modelos de decisión,
generación de código, recuperación aumentada, memoria a largo plazo, grafos de
conocimiento, simulación.
